\section{A linear strong-Pell multiplier in eleven integers}
\label{sec:linear-multiplier}

This section gives a separate improvement in the small-variable native
construction. It retains the complete source34 coefficient tuple, and hence
its inherited native/compiler/refutation calibration. The new polynomials
have eleven unrestricted integer unknowns. Their exact degree and live
binary-operation counts are
\[
  (440,126)\quad\hbox{and}\quad(412,127).
\]
The newest inspected paired predecessors have profiles $(456,124)$ and
$(416,125)$. The change therefore lowers degree while adding two operations.
It is a new tradeoff, with no assertion of a global minimum. All polynomial
abbreviations below are registers, not additional quantified unknowns.

The idea is to replace the square of the strong Pell multiplier by a linear
input. A coprime product forces that input to become a square at every zero.
Prime-adic orders of polynomial weights keep the radical terms separated, and a
rescaled interpolation request supplies the converse. The sign of the new
input cannot be assumed; recovering it is part of the proof.

\subsection{The literal polynomial}

Let
\[
 j,C,z,h,r,v,w,T,c,n,s\in\mathbb Z.
\]
The six fixed positive coefficients are
$\beta,K_0,M_C,M_F,A_c,B_c$. Their complete integer payloads and finite
reconstruction recipe are supplied with this component. In particular
\[
 \beta=2^{1064921}-1,\qquad A_c=2^{1113},\qquad B_c=2^{666104}.
\]
$M_F$ already includes its source shift. The other three large constants
are the exact whole source34 tuple, not new free parameters.
Define
\begin{align*}
 J&=6j+1,&q&=\beta J+1,&W&=12w+4,\\
 X&=qW,&Y&=qv,&E&=XY,\\
 a&=E+Y,&H&=4a+3,&\Delta&=a^2+H=(a+2)^2-1,\\
 F&=(K_0+W)C-z(q-1),&&&\\
 R&=\bigl(q(q-F)-C\bigr)(q^2-1)+(M_C+qM_F)J,&&&\\
 k&=R+1+hE,&\eta&=c-Yk.&&
\end{align*}
Set
\begin{align*}
 D_{\rm err}&=(A_cF-q)^2+(B_cC-q)^2,\\
 B_{\rm dg}&=k^2(q^2-D_{\rm err})-4q^2\eta^2,\\
 D&=3B_{\rm dg}-4,&\ell&=XY^2,&L&=\ell k,\\
 D_t&=1+L(L+k),&D_g&=1+DD_ts^2,&P&=D_tD_g.
\end{align*}
Here the new multiplier enters linearly:
\begin{align*}
 U&=rc^2,&I_2&=rc^4,&D_f&=1+\Delta I_2,\\
 Q&=\Delta^2I_2,&\delta&=Q-1,&V&=D_f(c-R+4UT)-c,\\
 C_y&=QV^2-1,&D_a&=\delta C_y,&K&=X(3c-2X)+2.
\end{align*}
Finally put
\begin{align*}
 M&=2q^2,&D_c&=MH,&S&=srM^5H,\\
 z_c&=D_cT+MK+H(v-q^2),&&&\\
 A_c^{\ast}&=(z_c^2+3D_c^2)(z_c^3-19D_c^3),&&&\\
 \mathcal A&=9S^2rD_f,&\mathcal B&=S^2D_a,
 &\mathcal G&=81S^2r^2P.
\end{align*}
The starred carrier $A_c^{\ast}$ is unrelated to the fixed cone coefficient
$A_c$. Use $Z=Sn-A_c^{\ast}$ for the direct version and
$Z=S(n+QV)-A_c^{\ast}$ for the centered version. The single equation is
\begin{equation}
 \boxed{\mathcal N=
 \left[(Z^2+\mathcal A-\mathcal B-\mathcal G)^2
       -4(\mathcal AZ^2+\mathcal B\mathcal G)\right]^2
       -64\mathcal AZ^2\mathcal B\mathcal G=0.}
 \label{eq:linear-multiplier-norm}
\end{equation}
The delivered circuits use only integer leaves and binary addition,
subtraction and multiplication. They absorb the fixed scale $81$ into
$D_g$ and reuse $S^2r$. The paired evaluation of the three-root norm is
credited to the incoming Round15 construction. Formally,
\[
 \mathcal N=\prod_{\epsilon\in\{-1,1\}^3}
 \bigl(Z-\epsilon_1\sqrt{\mathcal A}
        -\epsilon_2\sqrt{\mathcal B}
        -\epsilon_3\sqrt{\mathcal G}\bigr).
\]
This is a polynomial identity even when the radicands are dependent,
zero or negative; the radicals are used only in its proof.

\subsection{Signed bootstrap before Pell classification}

The authentic coefficient restrictions give
$q$ even, $q\equiv2\pmod3$, $|q|\ge16$, and $|X|\ge4|q|$.
If $v=0$, then $a=0$ and $\Delta=3$; otherwise $|a|>3$.
Thus $\Delta>1$ on every integer assignment. Also $H$ is odd and
nonzero, $H\equiv3\pmod q$, and $\gcd(H,M)=1$.
Since $X\equiv2\pmod3$, we have $a,H,\Delta\equiv0\pmod3$.
The consecutive product $\ell(\ell+1)$ is nonnegative, so $D_t\ge1$.

The polynomial $\pi(t)=(t^2+3)(t^3-19)$ has no rational root.
Consequently $A_c^{\ast}=D_c^5\pi(z_c/D_c)$ is never zero.
At $S=0$, equation~\eqref{eq:linear-multiplier-norm} reduces to
$(A_c^{\ast})^8=0$, impossible. Every zero therefore has $s,r\ne0$.

The auxiliary radicand $D_a$ is positive without assuming $r>0$.
If $c=0$, then $Q=0$ and $D_a=1$. If $c\ne0,r<0$, then $Q<0$ and
both factors in $(Q-1)(QV^2-1)$ are negative. If $c\ne0,r>0$, then
$Q>1$ and $D_f>|c|+1$. Because $V=D_fb-c$ for an integer $b$, either
$b=0$ and $|V|=|c|\ge1$, or $|V|\ge D_f-|c|>1$; again $D_a>0$.

The product $rD_f$ is nonzero and is positive whenever $c\ne0$.
Its only possible negative case is $c=0,r<0$, when $rD_f=r$.
After dividing a vanishing conjugate factor by $S$, its radical terms,
with independent signs, are
\begin{equation}
 3\sqrt{rD_f},\qquad \sqrt{D_a},\qquad 9r\sqrt P.
 \label{eq:linear-radicals}
\end{equation}
One negative radicand among $rD_f,P$ would leave a single nonzero
imaginary term. If both were negative, cancellation would require
$9(-r)=81r^2(-P)$, hence $1=9rP$, impossible over the integers.
Thus $rD_f>0$ and $P\ge0$. In fact $P>0$: the even fixed cone
coefficients make $B_{\rm dg}$ and $D$ divisible by $4$, so
$D_g\equiv1\pmod4$.

If $B_{\rm dg}\le0$, then $D_g\le-3$, contradicting $P>0$.
Therefore $B_{\rm dg}\ge4$. The usual centered disk inequality yields
$k\ne0$, the common strict sign of $q,F,C,q-F-C$, and
$|c/k-Y|<1/2$. Nonvanishing of $c$ is deferred until the cubic scale
has been recovered.

\subsection{A variable prime-adic separation lemma}

Write $a_3=v_3(|r|)$ and $b_3=v_3(P)$. The congruences
$D_f\equiv D_a\equiv1\pmod3$ imply
$v_3(rD_f)=a_3$ and $v_3(D_a)=0$.
Decompose the positive integer radicands in
\eqref{eq:linear-radicals} into squares times squarefree integers.
For a squareclass whose squarefree part has $3$-adic valuation
$\varepsilon\in\{0,1\}$, the valuations of its coefficients are,
when the corresponding terms occur,
\[
 \begin{array}{c|c}
 \text{auxiliary term}&0\quad(\varepsilon=0\text{ only})\\
 \text{strong term}&1+(a_3-\varepsilon)/2\\
 \text{guard term}&2+a_3+(b_3-\varepsilon)/2.
 \end{array}
\]
If the auxiliary term is present it has uniquely least valuation.
If strong and guard belong to the same class, their valuation difference
is $1+(a_3+b_3)/2>0$. Each nontrivial squareclass therefore has a unique
coefficient of least valuation and cannot cancel. Linear independence
of distinct squarefree radicals forces every radicand in
\eqref{eq:linear-radicals} to be an integer square.

In particular $rD_f$ is a positive square and $\gcd(r,D_f)=1$.
The two factors are either both positive squares or both negatives of
positive squares. The latter would give $-f^2\equiv1\pmod3$, impossible.
Thus
\begin{equation}
 r=u^2>0,\qquad D_f=f^2>0. \label{eq:linear-recovery}
\end{equation}
Also $\gcd(D_t,D_g)=1$ gives $D_t=\tau^2,D_g=g^2$.
The integral radical sum now implies $S\mid A_c^{\ast}$.

Since $M^5\mid S$, monicity of the homogenized quintic gives $M\mid z_c$:
if $p^e\Vert M$ but $v_p(z_c)=b<e$, the leading term $z_c^5$ would have
strictly smaller valuation than every other term of $A_c^{\ast}$.
Reducing modulo $M$ and using $\gcd(M,H)=1$ yields
\[
 v=q^2(2v_{\rm old}+1),\qquad Y=(2v_{\rm old}+1)q^3\ne0.
\]
If $c=0$, the positive budget would equal
$k^2[q^2(1-4Y^2)-D_{\rm err}]\le0$, a contradiction.

Now define the proof-only integer $I=c^2u$. We obtain
\[
 f^2-\Delta I^2=1,\quad Q=(\Delta I)^2>1,\quad
 V=f^2(c-R+4IT_{\rm old})-c,\quad T_{\rm old}=uT.
\]
The rational auxiliary root $\sqrt{D_a}/(Q-1)$ is integral by the
retained adjacent-coefficient Pell descent. Hence the old strong and
auxiliary interfaces, including $c^2\mid I$,
$V\equiv-c\pmod f$ and $V\equiv-R\pmod c$, are all restored.

\subsection{Source recovery and the converse}

The inherited signed rank argument applies to this recovered interface.
It requires neither oddness of $I/c^2$ nor freedom to change
$T_{\rm old}$ at this stage. It gives
\[
 c=\operatorname{sgn}(c)\psi_{a+2}(R),\quad
 R=2t+1,\ t\ge24,\quad
 |Y|>(|X|-3)^t/2-1,\quad |H|>|X|(|X|-3)^t.
\]
The divisibility $H\mid A_c^{\ast}$ gives $H\mid K^5$.
With $\sigma=\operatorname{sgn}(c)$ and
$\Pi=(\sigma X-2^R)(2^R\sigma X+1)$, the retained projection identity
therefore implies $H\mid\Pi^5$. The actual fixed $\beta$ gives
$x=|X|\ge2^{64}$, and
\[
 |\Pi|^5<2^{20t+20}x^{10}<x^{t+1}/2^t<|H|.
\]
The middle inequality follows from $64(t-9)>21t+20$ for $t\ge24$.
Thus $\Pi=0$ and $\sigma X=2^R$. The retained orientation and
centered half-binomial argument restores an actual old native source
zero, and its ordinary calibration yields a finite syntactic refutation.

Conversely, start with the old canonical outer witness supplied by a
finite refutation. It has $D_t=\tau^2$, $B_{\rm dg}\ge64$,
$D_t\equiv1\pmod3$, and $D\equiv2\pmod3$.
Choose a positive Pell solution $g^2-DD_ts^2=1$.
Choose one positive strong pair $I,f$ with $4sc^2\mid I$, and keep it
fixed. Put
\[
 u=I/c^2,\qquad r=u^2,\qquad m=sr=su^2.
\]
The inherited supported-modulus interpolation theorem supplies every
$T_{\rm old}$ residue modulo any positive modulus all of whose prime
divisors divide this fixed $I$. Its full proof is retained in
\texttt{source/POWERED\_PROOF.md}, Section~5.

Canonical projection gives $H\mid K$, so
$\kappa=K/H+(v-q^2)/M$ is integral. Choose an intersective root
$a_0\pmod m$ of $\pi(t+\kappa)$. Every prime divisor of $um=su^3$
divides $I$. Request
\[
 T_{\rm old}\equiv ua_0\pmod{um}.
\]
This yields an auxiliary Pell witness while retaining the same $I,f$.
Consequently $T=T_{\rm old}/u$ is integral and $T\equiv a_0\pmod m$.
Thus
\[
 \frac{A_c^{\ast}}S=\frac{H^4\pi(T+\kappa)}{sr}\in\mathbb Z.
\]
For the direct polynomial take
$n=A_c^{\ast}/S+3uf+\delta y+9r\tau g$; for the centered polynomial
subtract $QV$. One conjugate factor vanishes. The order of witness
choices is outer data, guard pair, fixed strong pair, residue,
auxiliary pair, and finally $n$. The strong pair is never changed after
the residue is requested.

Composing both directions with the retained calibration proves in
ordinary ZFC the equivalence between universal nonvanishing and
$\operatorname{Con}(\mathrm{ZFC})$. This is an ordinary proof composed
with named inherited theorems, not a new complete Lean, Rocq or encoded
object-language ZFC derivation. It proves neither side of that equivalence.

\subsection{Exact degrees and finite validation}

The relevant exact register degrees are
\[
 \begin{array}{c|rrrrrrrrrr}
 \text{register}&I_2&D_f&Q&V&D_a&D_t&P&S&A_c^{\ast}&
       (\mathcal A,\mathcal B,\mathcal G)\\ \hline
 \text{degree}&5&13&21&18&78&22&62&16&35&(46,110,96).
 \end{array}
\]
The direct version has uniquely dominant radicand $\mathcal B$, giving
degree $4\cdot110=440$. In the centered version let $C_0=SQV$ and
$Z_0=Sn-A_c^{\ast}$. The exact defect
\[
 \mathcal B-C_0^2=-S^2\bigl(Q(V^2+1)-1\bigr)
\]
gives $\deg(Z^2-\mathcal B)\le90$.
In the symmetric norm form the inner bracket has unique highest term
$-4\mathcal B\mathcal G$, of degree $206$; its other terms have
degrees at most $192,156,142$. The final correction has degree $362$.
Hence the exact centered degree is $412$.

Let $q_0=6\beta j$, $X_0=12q_0w$, $Y_0=q_0v$,
$R_0=-q_0^3(12wC-q_0z)$ and $k_0=R_0+hX_0Y_0$.
The entire centered highest form is
\[
 2^{64}3^{10}s^{12}r^{20}c^{32}q_0^{84}X_0^{40}Y_0^{52}R_0^4k_0^{12}.
\]
For example the following actual coefficients remain nonzero at the
complete fixed tuple:
\begin{align*}
 [j^{216}w^{52}v^{64}h^{12}z^4s^{12}r^{20}c^{32}]\mathcal N_{
 \rm centered}&=2^{384}3^{278}\beta^{216},\\
 [j^{224}w^{56}v^{56}z^8s^8r^{24}c^{64}]\mathcal N_{
 \rm direct}&=2^{392}3^{280}\beta^{224}.
\end{align*}
These establish exact total degree globally; they are not numerical
samples at particular solutions.

The separate finite checker imports neither the producer nor its DAG
utility. It authenticates the full coefficient tuple and retained source
files, checks all 68 register identities in each circuit over a free
integer polynomial ring, verifies every coefficient of the eight-factor
norm, checks the whole leading forms, and rejects seven changed nodes per
variant. Normal and optimized runs pass, the coefficient recipe reproduces
all six integers byte for byte, and both circuits rebuild identically.
These checks establish finite algebra and transcription claims; the
unbounded Pell and consistency interpretation remain the ordinary proofs
just described.
